Domestic electric storage water heaters may require users to be physically near the appliance to monitor its temperature and
operate its controls. 
The Rheem EW1-50RP (50 L, 1500 W) is the reference appliance for this project.
It measures and displays the temperature at the center of its tank, controls heating through an internal thermostat,
and provides an AI Smart operating mode~\cite{rheemmanual}.
However, access to its temperature readings and control interface is restricted to the appliance's physical panel.

This project addresses this limitation through an external Internet of Things (IoT) system that integrates sensing,
mechanical actuation, and wireless communication without modifying the heater's internal electrical circuitry.
An ESP32-CAM captures images of the original temperature display for optical recognition, while an RGB sensor monitors the heating indicator.
Additionally, a Hall-effect sensor supplies a separate pulse-acquisition channel for flow and accumulated-volume calculations.

An ESP32-WROOM microcontroller coordinates sensor acquisition, servo actuation, and communication through the MQTT protocol.
A Node.js application service processes camera images, manages MQTT messages, and communicates with a web-based dashboard
through WebSocket. This architecture enables remote visualization of temperature readings, sensor measurements, and device information,
while allowing control commands to be transmitted to the embedded controller.

The system is designed around the Rheem EW1-50RP while preserving the heater's original temperature regulation and
electrical protection mechanisms.
Its modular architecture also provides a reference for extending monitoring and external control capabilities to other
appliances with local-only interfaces.

